\documentclass[10pt,a4paper]{amsart}
\usepackage{amsmath,amssymb,amsthm,graphicx}
\usepackage{verbatim}
\newenvironment{reference}{\comment}{\endcomment}
%\newenvironment{reference}{}{}
\newenvironment{slogan}{\comment}{\endcomment}
\newenvironment{history}{\comment}{\endcomment}
% For commutative diagrams we use Xy-pic
\usepackage[all]{xy}

% We use 2cell for 2-commutative diagrams.
\xyoption{2cell}
% Theorem environments.
%
\theoremstyle{plain}
\newtheorem{theorem}[subsection]{Theorem}
\newtheorem{proposition}[subsection]{Proposition}
\newtheorem{lemma}[subsection]{Lemma}

\theoremstyle{definition}
\newtheorem{definition}[subsection]{Definition}
\newtheorem{example}[subsection]{Example}
\newtheorem{exercise}[subsection]{Exercise}
\newtheorem{situation}[subsection]{Situation}

\theoremstyle{remark}
\newtheorem{remark}[subsection]{Remark}
\newtheorem{remarks}[subsection]{Remarks}

\numberwithin{equation}{subsection}

% Macros
%
\def\lim{\mathop{\mathrm{lim}}\nolimits}
\def\colim{\mathop{\mathrm{colim}}\nolimits}
\def\Spec{\mathop{\mathrm{Spec}}}
\def\Hom{\mathop{\mathrm{Hom}}\nolimits}
\def\Ext{\mathop{\mathrm{Ext}}\nolimits}
\def\SheafHom{\mathop{\mathcal{H}\!\mathit{om}}\nolimits}
\def\SheafExt{\mathop{\mathcal{E}\!\mathit{xt}}\nolimits}
\def\Sch{\mathit{Sch}}
\def\Mor{\mathop{\mathrm{Mor}}\nolimits}
\def\Ob{\mathop{\mathrm{Ob}}\nolimits}
\def\Sh{\mathop{\mathit{Sh}}\nolimits}
\def\NL{\mathop{N\!L}\nolimits}
\def\CH{\mathop{\mathrm{CH}}\nolimits}
\def\proetale{{pro\text{-}\acute{e}tale}}
\def\etale{{\acute{e}tale}}
\def\QCoh{\mathit{QCoh}}
\def\Ker{\mathop{\mathrm{Ker}}}
\def\Im{\mathop{\mathrm{Im}}}
\def\Coker{\mathop{\mathrm{Coker}}}
\def\Coim{\mathop{\mathrm{Coim}}}

% Boxtimes
%
\DeclareMathSymbol{\boxtimes}{\mathbin}{AMSa}{"02}

%
% Macros for moduli stacks/spaces
%
\def\QCohstack{\mathcal{QC}\!\mathit{oh}}
\def\Cohstack{\mathcal{C}\!\mathit{oh}}
\def\Spacesstack{\mathcal{S}\!\mathit{paces}}
\def\Quotfunctor{\mathrm{Quot}}
\def\Hilbfunctor{\mathrm{Hilb}}
\def\Curvesstack{\mathcal{C}\!\mathit{urves}}
\def\Polarizedstack{\mathcal{P}\!\mathit{olarized}}
\def\Complexesstack{\mathcal{C}\!\mathit{omplexes}}
% \Pic is the operator that assigns to X its picard group, usage \Pic(X)
% \Picardstack_{X/B} denotes the Picard stack of X over B
% \Picardfunctor_{X/B} denotes the Picard functor of X over B
\def\Pic{\mathop{\mathrm{Pic}}\nolimits}
\def\Picardstack{\mathcal{P}\!\mathit{ic}}
\def\Picardfunctor{\mathrm{Pic}}
\def\Deformationcategory{\mathcal{D}\!\mathit{ef}}
\title{The kernel of a surjection onto a module of finite presentation is of finite type, locally on a ringed space}
\author{The Stacks Project Authors}
\date{stacks-project:modules.tex:local-snapshot-20261012; GNU FDL 1.2 or later}
\begin{document}
\maketitle
\noindent\textit{Editorial note.} All delivered mathematical text is the author's own, copied from the pinned Stacks Project chapter source sources/stacks-modules/pinned-original-source.tex; no statement or proof step was rewritten, reordered or removed.\par\medskip
\noindent\textit{Editorial note.} Preamble provenance: the Stacks Project keeps one single shared preamble for every chapter; this card takes it directly from the staged sources/stacks-modules/preamble.tex as three verbatim spans (lines 8-12, 14-18 and 154-217), so no derived file was written and nothing was added or redefined. The wrapper class and the base packages are supplied by the assembler.\par\medskip
\noindent\textit{Editorial note.} Bibliography provenance: neither the counted statement nor its proof contains a citation (this lane rescans the whole span before writing the card), so the card carries no bibliography and no entry was reconstructed.\par\medskip
\noindent\textit{Editorial note.} Selection provenance: this card was opened under the file-based control contract (CONTROL.md), with the mechanical pre-card gate run on its statement and proof interval before the card was written.\par\medskip
\noindent\textit{Allowed standard prerequisites (editorial).}\begin{itemize}
\item The notion of a module of finite presentation: locally there is a presentation by finite free modules. The author's section on modules of finite presentation, in which the counted statement occurs. Supplies the local presentation used throughout the proof.
\item The local lifting of sections of the quotient along a surjection and the generation of a kernel by displays. Used inside the counted proof at source lines 1698-1714. Produces the two maps whose combination surjects onto the kernel of the given surjection.
\end{itemize}\medskip
\begin{lemma}
\label{lemma-kernel-surjection-finite-free-onto-finite-presentation}
Let $(X, \mathcal{O}_X)$ be a ringed space.
Let $\mathcal{F}$ be an $\mathcal{O}_X$-module of finite presentation.
\begin{enumerate}
\item If $\psi : \mathcal{O}_X^{\oplus r} \to \mathcal{F}$ is a surjection,
then $\Ker(\psi)$ is of finite type.
\item If $\theta : \mathcal{G} \to \mathcal{F}$ is surjective with
$\mathcal{G}$ of finite type, then $\Ker(\theta)$ is of finite type.
\end{enumerate}
\end{lemma}

\begin{proof}
Proof of (1).
Let $x \in X$. Choose an open neighbourhood $U \subset X$ of $x$
such that there exists a presentation
$$
\mathcal{O}_U^{\oplus m}
\xrightarrow{\chi}
\mathcal{O}_U^{\oplus n}
\xrightarrow{\varphi}
\mathcal{F}|_U
\to
0.
$$
Let $e_k$ be the section generating the $k$th factor of
$\mathcal{O}_X^{\oplus r}$. For every $k = 1, \ldots, r$
we can, after shrinking $U$ to a small neighbourhood of $x$,
lift $\psi(e_k)$ to a section $\tilde e_k$ of
$\mathcal{O}_U^{\oplus n}$ over $U$. This gives a morphism
of sheaves $\alpha : \mathcal{O}_U^{\oplus r} \to \mathcal{O}_U^{\oplus n}$
such that $\varphi \circ \alpha = \psi$.
Similarly, after shrinking $U$, we can find a morphism
$\beta : \mathcal{O}_U^{\oplus n} \to \mathcal{O}_U^{\oplus r}$
such that $\psi \circ \beta = \varphi$. Then the map
$$
\mathcal{O}_U^{\oplus m} \oplus
\mathcal{O}_U^{\oplus r}
\xrightarrow{\beta \circ \chi, 1 - \beta \circ \alpha}
\mathcal{O}_U^{\oplus r}
$$
is a surjection onto the kernel of $\psi$.

\medskip\noindent
To prove (2) we may locally choose a surjection
$\eta : \mathcal{O}_X^{\oplus r} \to \mathcal{G}$.
By part (1) we see $\Ker(\theta \circ \eta)$ is of finite type.
Since $\Ker(\theta) = \eta(\Ker(\theta \circ \eta))$ we win.
\end{proof}

\end{document}
